% Ad Familiares 7.1 (ad-familiares-07-01) --- English translation
% Translated by Claude (Anthropic), from the Latin (Perseus).
% Style guide: STYLE.md.

\ciceroLetterOpener{M. CICERO S. D. M. MARIO.}

\ciceroSection{1} If some pain of the body, or some weakness of your health, has held you back from coming to the games, I assign it more to fortune than to your wisdom. But if those things which the rest wonder at, you have judged contemptible — and if, when by your health you could come, you yet did not wish to — at each I rejoice: both that you have been without bodily pain, and that you were strong of mind, when you neglected things which others wonder at without cause; provided only that the fruit of your leisure has stood you in good stead. Of which you could enjoy yourself wonderfully, since you were left almost alone in that lovely place. Nor do I doubt that, in that bedroom of yours from which you have pierced through and laid open the gulf of Stabiae, through those days you spent the morning hours in little readings, while those who left you in the country, half-asleep, were watching the public mimes. The rest of the parts of the day you spent in the delights you had prepared for yourself at your own pleasure. We, however, had to endure those things which Spurius Maecius had approved.

\ciceroSection{2} On the whole, if you ask, the games were of the most lavish equipment, but not to your taste — for I conjecture from my own. For first, by way of the honour, those had returned to the stage who, by way of the honour, I had supposed had withdrawn from the stage. As for your favourite, our Aesopus, he was such that everyone agreed he ought to retire. When he had begun to take an oath, his voice failed him at this place: ``\textit{si sciens fallo}'' (``if I knowingly deceive''). Why should I tell you of the rest? You know the other games. Which had not even that much grace which middling games are wont to have. For the spectacle of the equipment took away all the cheer; with which equipment, of course, I do not doubt that you have done without it with the most equal mind. For what delight is there in six hundred mules in \textit{Clytaemnestra}, or three thousand mixing-bowls in the \textit{Trojan Horse}, or various armour of foot and horse in some battle? Things which had popular wonder would have brought you no delight at all.

\ciceroSection{3} If during those days you gave attention to your Protogenes — provided that he read you anything rather than my speeches — surely you have had not a little more delight than any of us. For I do not think you missed the Greek or Oscan plays, especially since you can watch the Oscans even in your own town-council, and you so dislike the Greeks that you do not even go to your villa by the Greek way. Why should I think you missed the athletes, you who despised the gladiators? In which Pompey himself confesses that both his pains and his oil were lost. There are left the two beast-hunts in five days, magnificent, no one denies; but what delight can there be to a man of polish, when either a feeble man is torn by a most powerful beast, or a splendid beast is run through with the hunting-spear? Which yet, if they are to be seen, you have often seen; nor have we who watch these things seen anything new. The last day was of the elephants. There was great wonder of the crowd and of the throng, but no delight; rather, a kind of pity followed, and the impression that there is some kind of fellowship between that beast and the human race.

\ciceroSection{4} I, however, on these days of the stage games — lest I should seem to you not only \textit{blessed} but altogether free — almost burst myself in the trial of your friend Gallus Caninius. If I had as easygoing a people as Aesopus had, I should, by Hercules, gladly give up my craft, and live with you and with men like ourselves. For I was already weary even before, when both age and ambition urged me on, and when I was at last allowed not to defend whom I would not. But now indeed in this time there is no life. For I expect no fruit of my labour, and I am sometimes compelled to defend men who have not deserved best of me, at the asking of those who have deserved well.

\ciceroSection{5} So I look at last for every reason of living at my own discretion; and your manner of leisure I both vehemently praise and approve; and that you visit us less, this I bear with the more equal mind because, even if you were at Rome, neither I should be allowed to enjoy your charm, nor you mine (if there is any in me), on account of my most troublesome occupations. From which if I shall relax (for to be entirely free I do not demand), I shall surely teach you, who for many years have studied nothing else, what it is to live as a human being. Only do you keep up and protect that weakness of yours of health, as you do, that you may visit our villas and run about with me on a little litter together.

\ciceroSection{6} I have written you these things at greater length than I am wont, not from abundance of leisure but from love of you, because, by a certain letter of yours, you had — if you remember — half-invited me to write something of this kind, that you might less regret your missing the games. Which if I have attained, I rejoice; if not, this however I console myself with: that hereafter you will come to the games, and visit us, and you will not leave to my letters any hope of your delight.
